% !TEX root = ../enac_project_fr.tex
\chapter{Modélisation et mise en œuvre du positionnement GNSS}
\label{chap:modelisation}

Ce chapitre présente les modèles nécessaires au calcul de la position d'un
récepteur GNSS. La démarche part du SPP, fondé sur les mesures de code et les
données diffusées, puis introduit les observations et corrections nécessaires
au PPP. Les détails des algorithmes d'estimation sont donnés en
annexe~\ref{ann:estimation}.

\section{Du positionnement SPP au PPP}
\label{sec:theorie_positionnement}

\subsection{Positionnement SPP : observations et inconnues}

La pseudo-distance est obtenue en estimant le décalage temporel entre le code
reçu et une réplique du même code générée par le récepteur. Pour le récepteur
$r$, le satellite $s$ et la fréquence $i$, elle est modélisée en mètres par
\begin{align}
P_{r,i}^s &=
\rho_r^s+c(\delta t_r-\delta t^s)+T_r^s+I_{r,i}^s
+\delta_{\mathrm{Sagnac},r}^s+\delta_{\mathrm{Shapiro},r}^s
+\varepsilon_{P,r,i}^s.
\label{new:code}
\end{align}
avec :
\begin{itemize}
\item $P_{r,i}^s$ la pseudo-distance mesurée ;
\item $\rho_r^s=\|\mathbf r^s(t_e^s)-\mathbf r_r(t_r)\|$ la distance
géométrique entre les positions du satellite à l'émission et du récepteur
à la réception ;
\item $\delta t_r$ et $\delta t^s$ les biais d'horloge du récepteur et du
satellite, en secondes, et $c$ la vitesse de la lumière ;
\item $T_r^s$ et $I_{r,i}^s$ les retards troposphérique et ionosphérique ;
\item $\delta_{\mathrm{Sagnac},r}^s$ la correction liée à la rotation terrestre ;
\item $\delta_{\mathrm{Shapiro},r}^s$ le retard gravitationnel de propagation ;
\item $\varepsilon_{P,r,i}^s$ l'erreur résiduelle, comprenant notamment le bruit
de mesure et les effets non modélisés.
\end{itemize}

Les coordonnées sont exprimées dans un repère terrestre. Dans cette écriture,
les positions aux deux dates sont soustraites avant correction de la rotation
terrestre ; le terme de Sagnac corrige cette convention. Il ne doit pas être
ajouté une seconde fois si la rotation est déjà prise en compte dans le calcul
géométrique. Les biais instrumentaux de code sont traités selon la convention
des horloges, comme précisé ci-dessous.

Après correction des termes évaluables, l'observation s'écrit
\begin{align}
\underbrace{P_{r,i}^s+c\delta t^s-T_r^s-I_{r,i}^s
-\delta_{\mathrm{Sagnac},r}^s-\delta_{\mathrm{Shapiro},r}^s}
_{\text{Pseudo-distance corrigée }\overline P_{r,i}^s}
&=\rho_r^s+c\delta t_r+\varepsilon_{P,r,i}^s.
\label{new:corrected_code}
\end{align}
avec :
\begin{itemize}
\item $\overline P_{r,i}^s$ la pseudo-distance après application des corrections ;
\item les autres termes définis dans l'équation~\eqref{new:code}.
\end{itemize}

Pour une constellation, les inconnues à une époque sont les trois coordonnées
du récepteur et son biais d'horloge. Il faut au moins quatre observations
indépendantes et une géométrie permettant de distinguer ces quatre paramètres.
La présence de satellites supplémentaires apporte de la redondance.

\subsection{Corrections utilisées en SPP}
\label{new:corrections}

\subsubsection{Orbites, horloges et biais instrumentaux}

Les messages de navigation fournissent les éphémérides diffusées et les
paramètres de correction des horloges. Les premières permettent de calculer
la position du satellite à l'instant d'émission. Pour le GPS, la correction
d'horloge repose sur un polynôme auquel s'ajoute une correction relativiste
liée à l'excentricité de l'orbite~\cite{isgps200n,ashby2003} :
\begin{align}
\delta t^s(t)&=a_0^s+a_1^s(t-t_{\mathrm{oc}}^s)
+a_2^s(t-t_{\mathrm{oc}}^s)^2+\Delta t_{\mathrm{rel}}^s(t).
\label{new:clock}
\end{align}
avec :
\begin{itemize}
\item $a_0^s$, $a_1^s$ et $a_2^s$ les coefficients diffusés d'horloge ;
\item $t_{\mathrm{oc}}^s$ leur époque de référence ;
\item $\Delta t_{\mathrm{rel}}^s$ la correction relativiste périodique.
\end{itemize}

À cela s'ajoute, selon le signal utilisé, une correction de délai de groupe :
le TGD pour le GPS ou les BGD pour Galileo. Ces paramètres relient les biais
des signaux à la convention de l'horloge diffusée. Leur application dépend
donc de l'observable et de la combinaison de fréquences ; une correction
monofréquence ne se reporte pas telle quelle à une combinaison ionosphère-libre.

\subsubsection{Rotation terrestre et propagation relativiste}

La Terre tourne pendant le trajet du signal. La correction de Sagnac tient
compte de cette rotation entre l'émission et la réception. Le retard de
Shapiro décrit, quant à lui, l'allongement du temps de propagation dans le
champ gravitationnel terrestre~\cite{ashby2003}. Ces effets sont distincts de
la correction relativiste appliquée à l'horloge du satellite.

\subsubsection{Troposphère}

Le retard troposphérique dépend principalement de la pression, de la
température et de la teneur en vapeur d'eau de l'atmosphère. Il affecte le
code et la phase avec le même signe. Il est décomposé en contributions
hydrostatique et humide~\cite{saastamoinen1972} :
\begin{align}
ZTD&=ZHD+ZWD,\\
T_r^s&=M_{\mathrm{dry}}(E_r^s)ZHD+
M_{\mathrm{wet}}(E_r^s)ZWD.
\label{new:tropo}
\end{align}
avec :
\begin{itemize}
\item $ZTD$, $ZHD$ et $ZWD$ les retards zénithaux total, hydrostatique et humide,
en mètres ;
\item $E_r^s$ l'élévation du satellite ;
\item $M_{\mathrm{dry}}$ et $M_{\mathrm{wet}}$ les fonctions de projection
des retards zénithaux sur la ligne de visée.
\end{itemize}

Dans le traitement SPP considéré, la troposphère est corrigée par un modèle
de Saastamoinen alimenté par des données météorologiques conventionnelles.
Aucun paramètre troposphérique n'est estimé. L'estimation de la composante
humide sera introduite avec le PPP.

\subsubsection{Ionosphère : modèle de Klobuchar}

L'ionosphère contient des électrons libres qui modifient la propagation des
signaux. Cet effet est dispersif : au premier ordre, son amplitude varie
comme l'inverse du carré de la fréquence~\cite{hauschild2017observations}.
\begin{align}
I_{r,i}^s&=\frac{40{,}3}{f_i^2}\,\mathrm{STEC}_r^s.
\label{new:iono}
\end{align}
avec :
\begin{itemize}
\item $I_{r,i}^s$ le retard de code, en mètres ;
\item $f_i$ la fréquence du signal, en hertz ;
\item $\mathrm{STEC}_r^s$ le contenu électronique intégré le long du trajet,
en électrons par mètre carré ;
\item $40{,}3$ le coefficient du modèle du premier ordre dans ces unités.
\end{itemize}

Le modèle de Klobuchar fournit une estimation de ce retard à partir des
coefficients diffusés dans le message de navigation GPS, de la position
approchée du récepteur, de la direction du satellite et du temps
GPS~\cite{klobuchar1987}. Il représente l'ionosphère par une couche mince,
évalue le retard vertical au point traversé par le signal, puis le projette
sur la ligne de visée. Ses équations sont données en annexe~\ref{Klobuchar}.

Cette correction permet un traitement monofréquence, mais laisse subsister
une erreur ionosphérique résiduelle. Elle ne décrit pas toutes les variations
spatiales et temporelles du contenu électronique.

\subsubsection{Ionosphère : combinaison ionosphère-libre}

Deux fréquences permettent d'annuler le terme ionosphérique du premier ordre
sans l'estimer par un modèle. Pour le code, on forme~\cite{hauschild2017combinations}
\begin{align}
P_{r,\mathrm{IF}}^s&=\alpha P_{r,1}^s+\beta P_{r,2}^s,
&
\alpha&=\frac{f_1^2}{f_1^2-f_2^2},
&
\beta&=-\frac{f_2^2}{f_1^2-f_2^2}.
\label{new:if}
\end{align}
avec :
\begin{itemize}
\item $P_{r,\mathrm{IF}}^s$ la pseudo-distance combinée ionosphère-libre
(\textit{ionosphere-free}, IF) ;
\item $f_1$ et $f_2$ les deux fréquences utilisées ;
\item $\alpha$ et $\beta$ les coefficients vérifiant
$\alpha+\beta=1$ et $\alpha/f_1^2+\beta/f_2^2=0$.
\end{itemize}

La première identité conserve les termes indépendants de la fréquence ; la
seconde annule la contribution ionosphérique du premier ordre.
La combinaison amplifie toutefois le bruit des observations. Pour deux
erreurs non corrélées, sa variance vaut
\begin{align}
\sigma_{P,\mathrm{IF}}^2&=\alpha^2\sigma_{P,1}^2+\beta^2\sigma_{P,2}^2.
\label{new:if_variance}
\end{align}
avec :
\begin{itemize}
\item $\sigma_{P,1}^2$ et $\sigma_{P,2}^2$ les variances des mesures de code ;
\item $\sigma_{P,\mathrm{IF}}^2$ la variance de leur combinaison.
\end{itemize}

L'IF réduit donc l'erreur ionosphérique, mais ne garantit pas à elle seule une
dispersion plus faible que le traitement monofréquence. Les effets
ionosphériques d'ordre supérieur ne sont pas annulés.

\subsection{Estimation de la position SPP}

Les observations corrigées dépendent de façon non linéaire des coordonnées.
La résolution part d'une position approchée et ajuste les paramètres pour
réduire les écarts entre observations et valeurs calculées. Dans ce travail,
le SPP utilise les moindres carrés pondérés~\cite{verhagen2017least} :
\begin{align}
\widehat{\mathbf x}
&=\mathop{\mathrm{arg\,min}}_{\mathbf x}
[\mathbf y-\mathbf h(\mathbf x)]^{\mathsf T}\mathbf R^{-1}
[\mathbf y-\mathbf h(\mathbf x)],\\
\mathbf x&=[\,\mathbf r_r^{\mathsf T}\quad c\delta t_r\,]^{\mathsf T}.
\label{new:wls}
\end{align}
avec :
\begin{itemize}
\item $\mathbf y$ le vecteur des observations ;
\item $\mathbf h(\mathbf x)$ les observations calculées à partir du modèle ;
\item $\mathbf R$ la covariance des erreurs d'observation ;
\item $\widehat{\mathbf x}$ l'état estimé.
\end{itemize}

Une observation de variance élevée reçoit moins de poids. L'état est ajusté
par itérations, indépendamment à chaque époque. La linéarisation, les
équations de résolution et le critère d'arrêt sont détaillés en
annexe~\ref{ann:estimation}.

L'ajout de constellations augmente le nombre d'observations et peut améliorer
leur répartition géométrique. Il impose toutefois de tenir compte des
différences entre systèmes. Après mise en cohérence des échelles de temps,
un biais inter-systèmes résiduel (\textit{Inter-System Bias}, ISB) est estimé
pour chaque constellation supplémentaire. En prenant le GPS comme référence,
le terme d'horloge récepteur pour Galileo devient
\begin{align}
b_r^{\mathrm{GAL}}&=c\delta t_r+\mathrm{ISB}_r^{\mathrm{GAL}}.
\label{new:isb}
\end{align}
avec :
\begin{itemize}
\item $c\delta t_r$ le biais d'horloge commun référencé au GPS, en mètres ;
\item $\mathrm{ISB}_r^{\mathrm{GAL}}$ le biais additionnel Galileo, en mètres.
\end{itemize}

Cet ISB est ajouté au vecteur d'état. Le même principe s'applique aux autres
constellations : le nombre minimal de mesures dépend alors du nombre de
biais effectivement estimés et du rang du système.

\subsection{Passage au PPP : phase et produits précis}

La précision du SPP reste limitée par le bruit du code, les données
diffusées et les erreurs résiduelles des corrections. Le PPP exploite la
phase de porteuse et des produits précis pour réduire ces limitations.
Il ne nécessite pas de station de référence à proximité du
récepteur~\cite{zumberge1997,koubaHeroux2001}. La précision centimétrique visée
constitue également une base pour de futurs travaux de détermination précise
d'orbite (POD), non développés ici.

\subsubsection{Mesure de phase et ambiguïté}

Le récepteur compare la phase de la porteuse reçue à celle de son oscillateur.
Une fraction de cycle est mesurée avec une grande finesse. Après acquisition,
le récepteur suit également l'accumulation des cycles, mais le nombre entier
initial reste inconnu : c'est l'ambiguïté de phase.

La phase exprimée en mètres possède le modèle
\begin{align}
L_{r,i}^s&=\rho_r^s+c(\delta t_r-\delta t^s)+T_r^s-I_{r,i}^s
+B_{r,i}^s+\delta_{\mathrm{Sagnac},r}^s+
\delta_{\mathrm{Shapiro},r}^s+\varepsilon_{L,r,i}^s.
\label{new:phase}
\end{align}
avec :
\begin{itemize}
\item $L_{r,i}^s$ la mesure de phase exprimée en mètres ;
\item $B_{r,i}^s=\lambda_iN_{r,i}^s$ l'ambiguïté métrique, dans un modèle
où les biais instrumentaux de phase sont traités séparément ;
\item $\lambda_i=c/f_i$ la longueur d'onde et $N_{r,i}^s$ l'ambiguïté entière ;
\item $\varepsilon_{L,r,i}^s$ l'erreur résiduelle de phase ;
\item les autres termes définis dans l'équation~\eqref{new:code}.
\end{itemize}

L'ionosphère retarde le code mais produit une avance de phase, d'où son signe
opposé. Les corrections supplémentaires nécessaires au PPP sont introduites
dans la sous-section suivante.

L'ambiguïté reste constante sur un arc de suivi continu. Dans une solution
flottante, elle est estimée comme un paramètre réel. La résolution entière
constitue une étape supplémentaire, distincte du PPP flottant, qui exige
notamment des conventions de biais adaptées~\cite{teunissen2017ambiguity}.

\subsubsection{Orbites, horloges et biais précis}

Le PPP remplace les éphémérides et horloges diffusées par des produits précis
issus du traitement d'un réseau de stations. Les fichiers SP3 fournissent les
positions des satellites et les fichiers CLK leurs corrections d'horloge.
Ces données sont interpolées pour obtenir les valeurs à l'instant d'émission.

Les biais de signal doivent être cohérents avec les conventions de ces
horloges. Les produits BIA peuvent fournir des biais différentiels (DCB ou DSB)
ou des biais propres à une observable (OSB). Ils permettent de relier les
observations aux références des produits précis ; le TGD diffusé n'est donc
pas ajouté systématiquement. Il faut éviter de corriger deux fois un effet
déjà absorbé dans l'horloge ou dans les paramètres estimés.

\subsection{Modèle PPP retenu}

\subsubsection{Corrections supplémentaires}

La précision désormais visée impose de tenir compte d'effets secondaires
dans le bilan d'erreur du SPP :
\begin{itemize}
\item les décalages et variations des centres de phase des antennes (PCO et
PCV), décrits par les calibrations ANTEX ;
\item l'enroulement de phase (\textit{wind-up}), lié à l'orientation relative
des antennes et affectant la phase ;
\item les marées terrestres, qui déplacent la station sous l'action
gravitationnelle de la Lune et du Soleil.
\end{itemize}

Le PCO décrit le décalage entre un point de référence de l'antenne et son
centre de phase moyen. La PCV décrit la variation résiduelle avec la direction
du signal. Les corrections sont signées selon le modèle d'observation retenu.
Les déplacements de marée interviennent dans la position utilisée pour
calculer la distance géométrique.

\subsubsection{Estimation de la troposphère}

La composante hydrostatique est calculée par un modèle. La composante humide,
plus variable, est estimée conjointement avec la position. Deux gradients
peuvent compléter le modèle pour représenter une asymétrie azimutale :
\begin{align}
T_r^s&=M_{\mathrm{dry}}(E_r^s)ZHD+
M_{\mathrm{wet}}(E_r^s)ZWD
+\underbrace{M_{\mathrm{grad}}(E_r^s)
(G_N\cos A_r^s+G_E\sin A_r^s)}_{\text{terme dépendant de l'azimut}}.
\label{new:tropo_grad}
\end{align}
avec :
\begin{itemize}
\item $A_r^s$ l'azimut du satellite ;
\item $G_N$ et $G_E$ les paramètres de gradient nord et est ;
\item $M_{\mathrm{grad}}$ la fonction de projection associée, dont la
normalisation fixe les unités des paramètres de gradient ;
\item les autres termes définis dans l'équation~\eqref{new:tropo}.
\end{itemize}

Sans gradients, le dernier terme est supprimé. Les fonctions de projection
et leurs conventions sont détaillées en annexe~\ref{annexe:nmf}.

\subsubsection{Observations combinées et paramètres estimés}

L'IF est appliquée au code et à la phase, avec les coefficients de
l'équation~\eqref{new:if}. Après les corrections de biais et d'antenne, et de
wind-up pour la phase, le modèle s'écrit
\begin{align}
P_{r,\mathrm{IF}}^{s,\mathrm{corr}}&=
\mathcal G_r^s+\varepsilon_{P,r,\mathrm{IF}}^s,\\
L_{r,\mathrm{IF}}^{s,\mathrm{corr}}&=
\mathcal G_r^s+B_{r,\mathrm{IF}}^s+\varepsilon_{L,r,\mathrm{IF}}^s,\\
\mathcal G_r^s&=\rho_r^s+c(\delta t_r-\delta t_{\mathrm{prec}}^s)
+T_r^s+\delta_{\mathrm{Sagnac},r}^s+\delta_{\mathrm{Shapiro},r}^s,\\
B_{r,\mathrm{IF}}^s&=\alpha B_{r,1}^s+\beta B_{r,2}^s.
\label{new:ppp_model}
\end{align}
avec :
\begin{itemize}
\item $P_{r,\mathrm{IF}}^{s,\mathrm{corr}}$ et
$L_{r,\mathrm{IF}}^{s,\mathrm{corr}}$ les observations combinées corrigées ;
\item $\mathcal G_r^s$ le terme commun, indépendant de la fréquence ;
\item $\delta t_{\mathrm{prec}}^s$ l'horloge satellite issue du produit précis,
avec les corrections requises par sa convention ;
\item $\rho_r^s$ la distance calculée avec l'orbite précise et la position de
station corrigée des marées ;
\item $B_{r,\mathrm{IF}}^s$ l'ambiguïté IF en mètres, qui n'est pas directement
une ambiguïté entière ;
\item les $\varepsilon$ les erreurs résiduelles des observations combinées.
\end{itemize}

Pour une constellation et un modèle troposphérique avec gradients, l'état est
\begin{align}
\mathbf x&=[\,\mathbf r_r^{\mathsf T}\quad c\delta t_r\quad ZWD\quad
G_N\quad G_E\quad B_{r,\mathrm{IF}}^1\quad\cdots\quad
B_{r,\mathrm{IF}}^n\,]^{\mathsf T}.
\label{new:ppp_state}
\end{align}
avec :
\begin{itemize}
\item $\mathbf r_r$ les trois coordonnées de la station ;
\item $n$ le nombre d'arcs d'ambiguïté représentés dans l'état ;
\item les autres paramètres définis précédemment.
\end{itemize}
Les ISB sont ajoutés lorsque plusieurs constellations sont utilisées.

\subsection{Estimation séquentielle et continuité des observations}

Le PPP considéré utilise un filtre de Kalman étendu. À chaque époque, le
filtre prédit l'état et son incertitude à partir de l'époque précédente, puis
les met à jour avec les observations disponibles. Il conserve ainsi
l'information acquise au cours du suivi. Les équations sont données en
annexe~\ref{ann:estimation}.

Chaque paramètre doit recevoir une valeur initiale et une incertitude.
Un modèle d'évolution décrit également les variations admises entre deux
époques. La position d'une station statique, son horloge, la troposphère et
les ambiguïtés n'ont pas les mêmes propriétés temporelles et ne reçoivent
donc pas le même modèle.

La phase de convergence correspond à l'accumulation d'information permettant
de mieux séparer ces paramètres. Le code, bien que moins précis, reste utile
pour initialiser la position, l'horloge et les ambiguïtés.

Une perte de verrouillage est une interruption du suivi de la porteuse.
Elle peut provoquer un saut de cycle, c'est-à-dire une discontinuité du
comptage des cycles. L'ambiguïté antérieure n'est alors plus valable.
Le traitement contrôle les indicateurs du récepteur et des combinaisons
sensibles aux discontinuités. Lorsqu'un saut est détecté ou qu'une interruption
dépasse la durée admise, l'ambiguïté concernée est réinitialisée.
Les mécanismes détaillés sont conservés en annexe~\ref{ann:estimation}.

Le tableau suivant résume les deux traitements retenus dans ce travail.
Ces choix ne constituent pas une définition générale des estimateurs
utilisables en SPP ou en PPP.

\begin{table}[htbp]
\centering
\small
\begin{tabular}{|p{0.20\linewidth}|p{0.32\linewidth}|p{0.34\linewidth}|}
\hline
Élément & SPP retenu & PPP retenu \\\hline
Observations & Code & Code et phase \\\hline
Orbites et horloges & Diffusées & Produits précis \\\hline
Ionosphère & Klobuchar ou IF & IF \\\hline
Troposphère & Modèle a priori & Retard humide estimé, gradients optionnels \\\hline
Estimation & WLS par époque & EKF sur plusieurs époques \\\hline
Ambiguïtés & Absentes & Estimées sur les arcs continus \\\hline
\end{tabular}
\caption{Principales différences entre les traitements SPP et PPP retenus.}
\label{new:summary}
\end{table}

\section{Implémentation avec Orekit}
\label{sec:implementation_nouvelle}
% Partie réservée à la rédaction suivante.
